\begin{frame}{bag-of-words의 두 가지 문제}
  \begin{center}
    \includegraphics[width=0.92\linewidth]{ch14_bow_sparsity.png}
  \end{center}
  \vskip0.5em
  \begin{enumerate}
    \item \textbf{10만 차원, 99.99\%가 0인 희소 벡터}
    \item \textbf{단어 간 거리에 의미 없음}
  \end{enumerate}
\note{\begin{itemize}
\item 그 전에는 \kb{bag-of-words}가 표준 --- 어휘의 모든 단어를 나열해
  각 문장마다 \textbf{단어가 몇 번 등장하는지 세어서} 벡터로.
  (예: ``고양이는 귀엽다'' $\to$ $[1, 1, 0, \dots, 0]$)
\item \kb{차원이 어휘 크기와 같음} --- 10만 단어 어휘면 10만 차원,
      99.99\%가 0인 \textbf{극도로 희소}(sparse)한 벡터. \; 14.1의 PCA가
      이런 데이터에 효과적인 것도, ``1인 성분이 몇 개뿐''인 구조를
      알아서 압축해 주기 때문.
\item \kb{단어 간 ``거리''에 의미가 없음} --- one-hot에서는
      ``고양이''와 ``강아지''도 ``고양이''와 ``사과''와 \textbf{완전히
      같은 거리}(내적 0, 코사인 0). \; 의미를 담을 공간이 애초에
      설계되지 않은 것.
\end{itemize}}
\end{frame}
